\begin{textsample}{POS Dim 3 – rr – Score 11.00 – rr/rr\_inf\_9.txt}  \label{ex:f3_pos_158}
3.4.2 CORPO, \textbf{GESTOS} E MOVIMENTOS Na primeira infância, o corpo é o instrumento expressivo e comunicativo por excelência, que serve de suporte para o desenvolvimento emocional e mental, sendo essencial na construção de \textbf{afetos} e conhecimentos.
Esse campo destaca experiências nas quais o CORPO, os \textbf{GESTOS} e os MOVIMENTOS constituem linguagens das quais as \textbf{crianças}, desde cedo, fazem uso e que as orientam em relação ao mundo.
O corpo expressa e carrega consigo não somente características físicas e biológicas, mas também marcas de nosso pertencimento social que repercutem em quem somos e nas experiências que temos em relação ao gênero, etnia ou raça, classe, religião e à sexualidade.
O corpo revela nossa singularidade, nossa identidade pessoal e social.
Por meio do olhar, do tato, da audição, do paladar, do olfato, das sensações, da postura, da \textbf{mímica}, dos movimentos impulsivos ou coordenados, dos \textbf{gestos}, as \textbf{crianças}, desde \textbf{bebês}, exploram o mundo, estabelecem relações, expressam -se, \textbf{brincam} e produzem conhecimentos sobre si, sobre o outro, sobre o universo social e cultural.
Nesse contexto, as \textbf{crianças} \textbf{brincam} com seu corpo, se comunicam e se expressam por meio das diferentes linguagens, como música, dança, teatro, \textbf{brincadeiras} de faz de conta, no entrelaçamento entre corpo, emoção e linguagem.
Na Educação \textbf{Infantil}, o corpo das \textbf{crianças} e dos \textbf{bebês} ganha centralidade, pois ele é o partícipe privilegiado das práticas pedagógicas de \textbf{cuidado} físico, orientadas para a emancipação e a liberdade e não para a submissão.
As \textbf{crianças} conhecem e reconhecem com o corpo suas sensações, funções corporais e, nos seus \textbf{gestos} e movimentos, identificam as suas potencialidades e limites, desenvolvendo, ao mesmo tempo, a consciência sobre o que é seguro e o que pode ser um risco.
Também podem explorar e vivenciar um amplo repertório de movimentos, \textbf{gestos}, olhares, \textbf{sons} e \textbf{mímicas} com o corpo – individualmente ou em pares – descobrindo variados modos de ocupação e uso do espaço, como sentar com apoio, rastejar, engatinhar, escorregar, caminhar se \textbf{apoiando} em berços, mesas e cordas, saltar, escalar, equilibrar -se, correr, dar cambalhotas, alongar, ações sempre norteadas pelas interações e \textbf{brincadeiras}.
O referido campo destaca experiências ricas e diversificadas, em que \textbf{gestos}, \textbf{mímicas}, posturas e movimentos expressivos constituem uma linguagem vital com a qual as \textbf{crianças} percebem e expressam emoções, reconhecem sensações, interagem, \textbf{brincam}, ocupam espaços e neles se localizam, construindo conhecimento de si e do mundo.
Destaca -se também que a capacidade de nomear, identificar e ter consciência do próprio corpo, assim como a construção de uma autoimagem positiva, estão associadas às oportunidades oferecidas às \textbf{crianças} para expressão e conhecimento da cultura corporal da sociedade em que vivem.
A experiência que este campo propõe à \textbf{criança} vai além de movimentar partes do corpo ou deslocar -se no espaço, pois assume um importante papel para o desenvolvimento e a aprendizagem na infância, proporcionando o conhecimento de si e do mundo.

% matched lemmas: afeto, apoiar, bebê, brincadeira, brincar, criança, cuidado, gesto, infantil, mímico, som
\end{textsample}
